\section{Introduction}
\label{sec:introduction}

\subsection{Background: Kolmogorov-Arnold Networks as vector fields}

A standard Neural ODE \cite{chen2018neuralode} approximates an unknown
dynamical system $\dot{\mathbf{u}}(t) = \mathbf{F}(\mathbf{u}(t))$ by a
learned vector field $\mathbf{f}_\theta(\mathbf{u}(t))$, parameterized by a
multilayer perceptron: fixed nonlinear activations on each node, with all
learnable structure in the linear weights,
$\mathbf{y} = \sigma(\mathbf{W}\mathbf{x}+\mathbf{b})$. A Kolmogorov-Arnold
Network (KAN) \cite{liu2024kan} inverts this: every \emph{edge} carries its
own learnable univariate function $\phi(x)$, and each node simply sums its
incoming edges, $y_j = \sum_i \phi_{i,j}(x_i)$, following the
Kolmogorov-Arnold representation theorem -- any continuous multivariate
function on a bounded domain can be written as a finite composition of
continuous univariate functions and addition,
\begin{equation}
f(\mathbf{x}) = \sum_{q=1}^{2n+1}\Phi_q\!\left(\sum_{p=1}^{n}\phi_{q,p}(x_p)\right),
\label{eq:kolmogorov-arnold}
\end{equation}
so that no fixed nonlinearity needs to be chosen in advance -- the network
learns what shape each univariate function should take. In practice each
edge function is represented as a fixed smooth residual plus a learnable
expansion in a chosen basis,
\begin{equation}
\phi(x) = w_b\, b(x) + \sum_{i=1}^{G} c_i B_i(x), \qquad
b(x)=\mathrm{silu}(x) = \frac{x}{1+e^{-x}},
\label{eq:kan-edge}
\end{equation}
where $B_i(x)$ are the $G$ basis functions centred on a grid of $G$ points,
$c_i$ are their trainable coefficients, and the trainable weight $w_b$ sets
how much of the fixed residual $b(x)$ each edge keeps. Each edge therefore
has $G+1$ parameters. (The original KAN implementation uses $G+k$ B-spline
functions on $G$ intervals and a separate scale on the expansion; in the
formulation used here that scale is absorbed into the $c_i$.) The choice of $B_i$ is a
design decision, not a mathematical necessity: a Gaussian radial basis
function (RBF) grid, $B_i(x)=\exp(-(x-\mu_i)^2/2\sigma^2)$, gives each edge
smooth global support, while a degree-$k$ B-spline built by the Cox-de Boor
recursion gives each edge compact local support, zero outside its own knot
span. KAN-ODEs \cite{koenig2024kanode} put this construction inside a Neural
ODE's vector field: $\dot{\mathbf{u}}=\mathbf{f}_\theta(\mathbf{u})$ with
every edge of $\mathbf{f}_\theta$ built from Equation~\ref{eq:kan-edge}, and
the resulting trajectory obtained by numerically integrating that field
forward in time.

\subsection{Architecture, and what was built rather than imported}

The reference KAN-ODE of this report, the base paper's Lotka-Volterra model,
is a two-layer network with widths $[2,10,2]$, each edge carrying the
expansion of Equation~\ref{eq:kan-edge} over a grid of $G=5$ basis
functions, wrapped inside an explicit Runge-Kutta integrator that advances
the state and hands the resulting trajectory to the loss. The other systems
use the same structure at a different size (Table~\ref{tab:train-settings}).
Figure~\ref{fig:kan-ode-architecture} shows the structure schematically, and
marks the three places where this report changes it.

\begin{figure}[htbp]
\centering
\resizebox{\textwidth}{!}{%
\begin{tikzpicture}[
  font=\small,
  nd/.style={circle, draw=ink, fill=white, minimum size=7mm, inner sep=0pt},
  edgefn/.style={-{Stealth[length=1.2mm]}, draw=ink!30, thin},
  blocklbl/.style={font=\small\sffamily\bfseries\color{ink}},
  ablationmark/.style={draw=accent, fill=white, circle, minimum size=5.5mm, inner sep=0pt, font=\bfseries\small\color{accent}}
]
  % small "basis icon" glyph standing in for one edge's learned function
  \newcommand{\basisicon}[3]{%
    \node[draw=ink!45, fill=white, minimum width=5mm, minimum height=3.6mm, inner sep=0pt] at (#1,#2) {%
      \begin{tikzpicture}[scale=0.15]
        \draw[ink!75, thick] plot[smooth] coordinates {(-1,#3) (-0.3,-0.6) (0.4,0.7) (1,-0.2)};
      \end{tikzpicture}};
  }
  % ===== outer wrapper: "KAN-ODEs" (background layer, drawn first so it
  % never occludes the block contents drawn on top of it) =====
  \begin{pgfonlayer}{background}
    \node[draw=green!30!black, thick, rounded corners=8pt, fill=green!4,
          fit={(-2.35,-3.0) (9.35,3.3)},
          label={[font=\normalsize\sffamily\bfseries\color{green!30!black}]above:%
            KAN-ODEs: Kolmogorov-Arnold Network Ordinary Differential Equations}] (outer) {};
  \end{pgfonlayer}

  % ===== Block 1: Kolmogorov-Arnold Network (spans x in [-2.0,2.0]) =====
  \node[nd] (u1) at (-0.8,-1.9) {$u_1$};
  \node[nd] (u2) at (0.8,-1.9) {$u_2$};
  \foreach \i/\x in {1/-1.5,2/-0.5,3/0.5,4/1.5} {
    \node[nd, minimum size=6mm] (h\i) at (\x,0) {\tiny$h_{\i}$};
  }
  \node[nd] (o1) at (-0.8,1.9) {\scriptsize$\dot u_1$};
  \node[nd] (o2) at (0.8,1.9) {\scriptsize$\dot u_2$};
  \foreach \i in {1,...,4} {
    \draw[edgefn] (u1) -- (h\i);
    \draw[edgefn] (u2) -- (h\i);
    \draw[edgefn] (h\i) -- (o1);
    \draw[edgefn] (h\i) -- (o2);
  }
  \foreach \x/\c in {-1.5/0.4,-0.5/-0.5,0.5/0.5,1.5/-0.3} { \basisicon{\x}{-0.95}{\c} }
  \foreach \x/\c in {-1.5/-0.4,-0.5/0.5,0.5/-0.5,1.5/0.3} { \basisicon{\x}{0.95}{\c} }
  \node[font=\tiny\sffamily\color{ink!70}, anchor=east] at (-2.1,1.9) {Output};
  \node[font=\tiny\sffamily\color{ink!70}, anchor=east] at (-2.1,0) {Hidden};
  \node[font=\tiny\sffamily\color{ink!70}, anchor=east] at (-2.1,-1.9) {Input};
  \begin{pgfonlayer}{background}
    \node[draw=green!35!black, thick, rounded corners=4pt, fill=green!10, fit={(-2.0,-2.55) (2.0,2.55)}, label={[blocklbl]above:Kolmogorov-Arnold Network}] (b1) {};
  \end{pgfonlayer}
  \node[ablationmark] at ([xshift=-6pt,yshift=-6pt]b1.north east) {1};

  % ===== Block 2: Ordinary Differential Equation (forward solve) =====
  \begin{scope}[xshift=7.0cm]
    \node[align=center, font=\scriptsize] (eqn) at (0,1.6)
      {$\mathbf{u}(t)=\mathbf{u}_0+\displaystyle\int_{t_0}^{t}\mathbf{f}_\theta(\mathbf{u}(\tau))\,\dd\tau$};
    \draw[thick, accent] plot [smooth] coordinates {(-1.5,0.05) (-0.5,-0.45) (0.5,0.2) (1.5,-0.2)};
    \foreach \px/\py/\lbl in {-1.5/0.05/u(0), -0.5/-0.45/u(t_i), 0.5/0.2/u(t_{i+1}), 1.5/-0.2/u(t_f)} {
      \fill (\px,\py) circle (1.2pt);
      \node[font=\tiny] at (\px,\py+0.35) {$\lbl$};
    }
    \draw[ink!55, thick] (-1.8,-1.1) -- (1.8,-1.1);
    \foreach \px/\tlbl in {-1.5/0, -0.5/t_i, 0.5/t_{i+1}, 1.5/t_f} {
      \draw[ink!55, thick] (\px,-1.04) -- (\px,-1.16);
      \node[font=\tiny] at (\px,-1.4) {$\tlbl$};
    }
    \begin{pgfonlayer}{background}
      \node[draw=green!35!black, thick, rounded corners=4pt, fill=green!10, fit={(-2.0,-1.7) (2.0,2.0)}, label={[blocklbl]above:Ordinary Differential Equation}] (b2) {};
    \end{pgfonlayer}
    \node[ablationmark] at ([xshift=-6pt,yshift=-6pt]b2.north east) {2};
    \node[font=\tiny\sffamily\color{ink!60}] at (0,-1.9) {Forward solve};
  \end{scope}


  % ===== two-way information exchange between the network and the solver:
  % the KAN network outputs the field du/dt, which flows into the solver;
  % the solver's running state u flows back into the network to evaluate
  % the field at the next step =====
  \draw[-{Stealth[length=2.4mm]}, draw=accent!85, thick]
    (b1.10) .. controls +(1.2,0.6) and +(-1.2,0.6) .. (b2.170)
    node[midway, above, font=\tiny\sffamily\color{accent!85}] {Gradient information ($\dd\mathbf{u}/\dd t$)};
  \draw[-{Stealth[length=2.4mm]}, draw=ink!65, thick]
    (b2.190) .. controls +(-1.2,-0.6) and +(1.2,-0.6) .. (b1.-10)
    node[midway, below, font=\tiny\sffamily\color{ink!65}] {State information ($\mathbf{u}=[x,y]$)};

  % ===== Block 3: Adjoint Sensitivity Analysis (backward solve), sized to
  % match the outer wrapper's height so the return arrow lands symmetrically =====
  \begin{scope}[xshift=13.2cm]
    \node[align=center, font=\scriptsize, text width=3.6cm] at (0,2.6)
      {gradient computation: direct autograd \emph{or} adjoint sensitivity};
    % the augmented adjoint system actually solved backward in time,
    % consistent with the state/adjoint/parameter-gradient triple of Eq. 12
    \node[align=center, font=\tiny] at (0,1.5)
      {$\dfrac{\dd}{\dd t}\!\begin{bmatrix}\mathbf{u}\\ \mathbf{a}\\ \partial L/\partial\theta\end{bmatrix}
        = \begin{bmatrix}\mathbf{f}_\theta(\mathbf{u})\\[2pt]
        -\mathbf{a}^\top\,\partial\mathbf{f}_\theta/\partial\mathbf{u}\\[2pt]
        -\mathbf{a}^\top\,\partial\mathbf{f}_\theta/\partial\theta\end{bmatrix}$};
    \draw[thick, ink!70] plot [smooth] coordinates {(1.6,-0.65) (0.6,-1.15) (-0.4,-0.5) (-1.6,-0.9)};
    \foreach \px/\py in {1.6/-0.65, 0.6/-1.15, -0.4/-0.5, -1.6/-0.9} {
      \fill (\px,\py) circle (1.2pt);
    }
    \node[font=\tiny] at (1.6,-0.3) {$\mathbf{a}(t_1)$};
    \node[font=\tiny] at (-1.6,-0.55) {$\mathbf{a}(t_0)$};
    \draw[{Stealth[length=2mm]}-, ink!55, thick] (-1.8,-1.8) -- (1.8,-1.8);
    \node[font=\tiny] at (1.6,-2.1) {$t_1$};
    \node[font=\tiny] at (-1.6,-2.1) {$t_0$};
    \node[font=\tiny\sffamily\color{ink!60}] at (0,-2.5) {Backward solve};
    \begin{pgfonlayer}{background}
      \node[draw=blue!40!black, thick, rounded corners=4pt, fill=blue!6, fit={(-2.0,-3.0) (2.0,3.3)}, label={[blocklbl]above:Adjoint Sensitivity Analysis}] (b3) {};
    \end{pgfonlayer}
    \node[ablationmark] at ([xshift=-6pt,yshift=-6pt]b3.north east) {3};
  \end{scope}

  % ===== single consolidated ablation legend, top right of the whole figure =====
  \node[draw=accent, thick, rounded corners=3pt, fill=white, align=left,
        font=\tiny\sffamily, inner sep=5pt,
        anchor=south east] at (15.2,5.3)
    {\textbf{\color{accent}Ablated in this report:}\\
     (1) edge basis -- Kolmogorov-Arnold Network\\
     (2) integrator -- Ordinary Differential Equation\\
     (3) differentiation method -- Adjoint Sensitivity Analysis};

  % ===== forward connector: outer wrapper to the adjoint box =====
  \draw[-{Stealth[length=3mm]}, draw=ink!65, line width=1.6pt]
    (outer.east) -- node[above, font=\tiny\sffamily\color{ink!70}]{Solution} (b3.west);

  % ===== training-loop back-edge: "KAN model update via gradient descent" =====
  \draw[-{Stealth[length=2.8mm]}, draw=blue!45!black, ultra thick]
    (b3.south) .. controls +(0,-1.0) and +(0,-1.0) .. (outer.south)
    node[midway, below=3pt, font=\tiny\sffamily\color{blue!45!black}] {KAN model update via gradient descent};
\end{tikzpicture}%
}
\caption{The KAN-ODE training loop, redrawn in the style of the base paper's
own architecture figure \cite{koenig2024kanode}. The legend at top right
lists this report's three ablated components and which block of the diagram
each corresponds to.}
\label{fig:kan-ode-architecture}
\end{figure}

With $G=5$ RBF centers per edge, layer one ($2\to10$) contributes
$10\times(2\times5+2)=120$ parameters and layer two ($10\to2$) contributes
$2\times(10\times5+10)=120$ parameters, for the $240$ total the base paper
reports. The reference Python KAN library, \texttt{pykan} \cite{liu2024kan},
represents every edge with its own B-spline object and refits its grid
during training. The base paper does not use it: its Julia code
\cite{koenig2024kanoderepo} builds the network from the \texttt{KDense} layer
of \texttt{Kolmogorov\-Arnold.jl} \cite{puri2024kolmogorovarnoldjl}, which
evaluates every edge's basis functions in one array operation followed by a
single matrix product. This project ports that layer to PyTorch, which makes
every basis substituted into Equation~\ref{eq:kan-edge} a one-function change
evaluated by an identical code path
(Section~\ref{sec:numerics}, Algorithm~\ref{alg:kdense}). The explicit
Runge-Kutta integrators, the training loop and all diagnostics were written
for this project rather than imported, so every basis and solver could be
instrumented identically. The base paper runs its integration loop in
Julia's SciML ecosystem, while this project runs it in PyTorch, so wall-clock
figures reported later are comparable across this project's own runs but
not directly against the paper's reported times.

\subsection{What the base paper fixed, and what this project varies}

The base paper reports two central claims on the Lotka-Volterra
predator-prey system: a KAN-ODE with 240 parameters extrapolates
Lotka-Volterra's limit cycle accurately from a short training window, and it
reaches a target training loss in roughly $10^4$ epochs against a
parameter-matched multilayer-perceptron Neural ODE's $10^5$ -- a claimed
tenfold training-speed advantage at 252 parameters. To obtain these results,
the paper commits to three specific design decisions that this project
treats as free variables rather than fixed choices: the \textbf{integrator}
(an adaptive fifth-order Runge-Kutta method, Tsit5), the \textbf{basis}
(Gaussian RBF), and the \textbf{differentiation method} used to compute
gradients through the integrator (continuous adjoint sensitivity, motivated
by its constant memory cost regardless of trajectory length).

The basis choice is not a detail invented for this project to vary: the
original Kolmogorov-Arnold Network paper \cite{liu2024kan} used the cubic
B-spline as its default edge basis, and the base paper replaces it with a
Gaussian RBF grid for the ODE setting. The two differ in ways that matter
once the edge function sits inside a differentiated integrator: the Gaussian
is infinitely differentiable, nonzero everywhere and a single closed-form
expression, while the cubic B-spline is only twice continuously
differentiable, exactly zero outside a few grid cells, and built by a
recursion that is several times more expensive to evaluate
(Section~\ref{sec:numerics}). The switch is a trade-off, not an obvious
improvement, so Section~\ref{sec:phase2} restores the B-spline and adds five
further bases (Chebyshev, Lagrange, Newton and two rational kernels) to
measure it directly rather than assume it.

Each of those three fixed choices maps directly onto one of this report's
studies: Section~\ref{sec:track-a} tests whether the adjoint method's
memory motivation actually holds at this project's scale; Sections
\ref{sec:phase2}, \ref{sec:track-b}, and \ref{sec:track-d} test the
integrator choice, sweeping solver order and, in the damping study (originally intended as a stiffness probe),
solver order jointly with the pendulum's damping coefficient; and
Section~\ref{sec:track-c} tests the basis choice directly, by making the
basis itself a learnable, trained quantity rather than a fixed decision made
in advance. Section~\ref{sec:track-e} extends the comparison outside the
base paper's own scope entirely, benchmarking KAN-ODE against a
sparse-regression alternative built for the same symbolic-discovery goal the
RBF basis was partly chosen to support, on both synthetic Lotka-Volterra data and a
synthetic, non-mechanistic epidemic curve.

\subsection{Why training itself needs to be treated as an object of study}

This project uses two testbeds beyond Lotka-Volterra -- the SIR epidemic
model and a damped pendulum -- chosen for simple, well-understood governing
equations, so any training failure is attributable to the architecture or
optimization rather than to ambiguous target dynamics. The integrator's
interaction with training is not predictable from classical numerical-analysis
theory alone, because it sits inside a backpropagation graph rather than a
standalone forward solve: a solver that is stable at the step size used
can still train badly, train to a wrong answer, or fail to generalize.
Diagnosing failures like these means treating training itself as an object of
study, not just the final trajectory error.

\subsection{Structure of this report}

Three method sections come first: the datasets, splits and training settings
(Section~\ref{sec:setup}); the numerical implementation of the KAN layer,
the integrators, the training loop and the structural fixes, with their
algorithms and cost (Section~\ref{sec:numerics}); and the verification of
that implementation against known solutions and an automated test suite
(Section~\ref{sec:verification}). The results then follow the order the work
was carried out. A baseline recipe is reproduced and used to establish the
evaluation protocol (Section~\ref{sec:phase1}), then stress-tested with a solver and
basis-function ablation that surfaces and fixes two training failures
(Section~\ref{sec:phase2}). Five further studies build on that fixed
baseline: gradient-method cost (Section~\ref{sec:track-a}), whether gradient
noise explains a solver-order effect (Section~\ref{sec:track-b}), a learnable
basis blend (Section~\ref{sec:track-c}), solver-damping interaction (planned as a stiffness test)
(Section~\ref{sec:track-d}), and KAN-ODE against sparse regression on
synthetic data and a non-mechanistic epidemic curve
(Section~\ref{sec:track-e}). A validation stage (Section~\ref{sec:phase4})
retrains key configurations at up to five times the epoch budget and audits
this project's recipe, and the base paper's claims, against the paper
itself.
Section~\ref{sec:discussion} draws together cross-cutting patterns, and
Section~\ref{sec:conclusion} closes the report.
